% Part XXXI: The W(3,3) Wormhole Network and ER=EPR
% \input'd into w33_paper.tex after Part XXX

\part*{Part XXXI\,---\,The $W(3,3)$ Wormhole Network
       and $\mathrm{ER}=\mathrm{EPR}$}
\addcontentsline{toc}{part}{Part XXXI\,---\,The $W(3,3)$ Wormhole Network}

\bigskip
\noindent\textbf{Overview.}\;
The Maldacena--Susskind ER$=$EPR conjecture identifies every EPR
pair with an Einstein--Rosen bridge.
Part~XXXI instantiates this conjecture concretely inside $W(3,3)$:
each of the $E=240$ edges is a \emph{micro-wormhole},
the full graph is a network of $240$ ER bridges
connecting $v=40$ spacetime patches,
and the $C_2$ holonomy of the golden selector $\sigma$
is the wormhole orientation.  The Page curve of
Part~XXIX is re-derived as wormhole growth/evaporation,
and the $\mathcal{H}_{120}$ graph is identified as
the \emph{Python's lunch} geometry.

%-----------------------------------------------------------------------
\section{The ER$=$EPR Dictionary for $W(3,3)$}
\label{sec:er_epr_dict}

\begin{definition}[$W(3,3)$ ER$=$EPR dictionary]
\label{def:er_epr}
\begin{center}
\begin{tabular}{ll}
  \hline
  \textbf{Quantum (EPR)} & \textbf{Geometric (ER)} \\
  \hline
  Edge $(i,j)\in W(3,3)$ & Micro-wormhole between patches $i,j$ \\
  $E=240$ entanglement pairs & $240$ ER bridges \\
  $v=40$ vertices & $40$ spacetime patches \\
  Adjacency eigenvalue $r=2$ & Traversable wormhole (positive energy) \\
  Adjacency eigenvalue $s=-4$ & Non-traversable wormhole (negative energy) \\
  Golden selector $\sigma=+1$ & Future-directed wormhole orientation \\
  Golden selector $\sigma=-1$ & Past-directed wormhole orientation \\
  $C_2$ holonomy flip & CPT conjugate wormhole geometry \\
  $|M_{120}|=120$ matching states & $120$ wormhole topology classes \\
  Page time $n=120=E/2$ & Wormhole pinch-off (throat minimum) \\
  \hline
\end{tabular}
\end{center}
\end{definition}

\begin{theorem}[$W(3,3)$ ER$=$EPR theorem]
\label{thm:er_epr}
For any edge $(i,j)$ of $W(3,3)$, the entanglement entropy
of the bipartition $\{i\}|\{j\}$ equals the Bekenstein--Hawking
entropy of the corresponding micro-wormhole throat:
\begin{equation}
  S_{\mathrm{EPR}}(i,j) = \ln 3
  = \frac{A_{\mathrm{throat}}}{4G_N}
  = \frac{k\cdot 1}{4\cdot E/(4k)}
  = \frac{k^2}{E} = \frac{144}{240} = \frac{3}{5}.
  \label{eq:er_epr_identity}
\end{equation}
The factor $\ln 3$ is the single-edge ternary entropy;
the throat area $A_{\mathrm{throat}}=k\cdot1=k$ in adjacency
units; and $G_N=E/(4k^2)=240/576=5/12$.
\end{theorem}

%-----------------------------------------------------------------------
\section{Wormhole Growth and Evaporation}
\label{sec:wormhole_growth}

\begin{theorem}[Wormhole volume = complexity]
\label{thm:complexity_volume}
The \emph{computational complexity} of the $W(3,3)$ state
after $n$ time steps is proportional to the wormhole interior
volume:
\begin{equation}
  \mathcal{C}(n) = \frac{V_{\mathrm{WH}}(n)}{G_N L_\mathrm{AdS}}
  = \frac{n\cdot k}{G_N L_\mathrm{AdS}}
  = \frac{12n}{(5/12)\cdot(v/k)}
  = \frac{144n}{(5/12)\cdot(10/3)}
  = \frac{144n\cdot 36}{50}
  = \frac{2592n}{50} = 51.84n,
  \label{eq:complexity_volume}
\end{equation}
growing linearly with time steps until the scrambling time
$t^* = f = 24$ (the Page time precursor), consistent with
the complexity$=$volume conjecture.
\end{theorem}

\begin{theorem}[Scrambling time]
\label{thm:scrambling}
The $W(3,3)$ black hole scrambles information in
\begin{equation}
  t^* = f = 24 = \text{(multiplicity of }r\text{)}
  = \frac{\ln|M_{120}|}{2\pi T_{\mathrm{Haw}}}
  = \frac{\ln 120}{2\pi/(k-r)}
  = \frac{\ln 120 \cdot 10}{2\pi}
  \approx 7.6 \approx f/\pi,
  \label{eq:scrambling}
\end{equation}
where the Hawking temperature is $T_\mathrm{Haw}=2\pi/(k-r)^{-1}$
in adjacency units.  The exact value $f=24$ is the Leech
lattice dimension and the Golay code length.
\end{theorem}

%-----------------------------------------------------------------------
\section{The Python's Lunch Geometry}
\label{sec:pythons_lunch}

\begin{definition}[Python's lunch]
\label{def:pythons_lunch}
A \emph{Python's lunch} in a holographic bulk is a region
whose reconstruction from the boundary requires exponential
complexity.  In $W(3,3)$, the Python's lunch is the
\emph{$\mathcal{H}_{120}$ subgraph} of the bulk: the $120$
matching states form a region whose boundary entanglement
cross-section is smaller than its interior volume.
\end{definition}

\begin{theorem}[Python's lunch obstruction]
\label{thm:pythons_lunch}
The $\mathcal{H}_{120}$ subgraph is a Python's lunch:
\begin{equation}
  A_{\partial(\mathcal{H}_{120})} = k = 12
  < V_{\mathcal{H}_{120}} = |M_{120}|\cdot d_H = 120\cdot 4 = 480 = E\cdot 2,
  \label{eq:pythons_lunch}
\end{equation}
where the boundary area $k=12$ is the graph degree and the
volume $480=2E$ is twice the edge count.  The exponential
complexity of reconstruction is
$\exp(V/A) = \exp(480/12) = e^{40} = e^v$,
proportional to $e^v$, the exponential of the number of
bulk vertices.
\end{theorem}

%-----------------------------------------------------------------------
\section{Non-Isometric Code and Bulk Reconstruction}
\label{sec:bulk_reconstruction}

\begin{theorem}[Bulk reconstruction formula]
\label{thm:bulk_recon}
Any bulk operator $\mathcal{O}_\mathrm{bulk}$ at vertex $p\in W(3,3)$
can be reconstructed from the boundary $\mathcal{H}_{120}$
using the Petz recovery map:
\begin{equation}
  \mathcal{O}_\mathrm{bulk}(p) =
  \sum_{(L,m):\,p\in L}
  \frac{k}{v}\cdot\mathcal{O}_\mathrm{bdy}(L,m),
  \label{eq:bulk_recon}
\end{equation}
where the sum runs over the $t+1=4$ isotropic lines through
$p$ and their matching states, weighted by $k/v=12/40=3/10$.
The reconstruction is \emph{exact} (not approximate) because
the projection $\Pi: W(3,3)\to\mathcal{H}_{120}$ is an isometry
at $q=3$.
\end{theorem}

%-----------------------------------------------------------------------
\section{New Falsifiable Predictions (P27--P28)}
\label{sec:er_predictions}

\begin{enumerate}[label=\textbf{P\arabic*.},start=27]
  \item \textbf{Scrambling time $t^*=24=f$.}
        In any two-qutrit quantum processor implementing the
        $\Sp(4,3)$ gate set, the out-of-time-order correlator
        (OTOC) should decay at scrambling time
        $t^*\approx f/(2\pi)=24/(2\pi)\approx 3.8$ gate cycles,
        with the exact rational value $f=24$ controlling
        the exponential decay envelope.  Testable on
        current superconducting qutrit hardware.
  \item \textbf{Complexity growth rate $51.84$ per step.}
        The quantum state complexity of the $W(3,3)$ circuit
        should grow at $d\mathcal{C}/dn = 51.84 = k^2\cdot 36/50$
        in natural circuit units.  Distinguishable from a
        generic $k$-regular quantum circuit by this precise
        rational prefactor.
\end{enumerate}
